\section*{Exact FEI and upper-set-potential audit of point-defect towers}

All Fourier transforms use the uniform probability measure, all logarithms
are base two, and zero-mass summands are zero.  For a spectral law
$p_S=\widehat f(S)^2$, put
\[
 \pi_f(S)=\sum_{T\supseteq S}p_T,
 \qquad
 U(f)=\sum_S p_S\log_2\frac{\pi_f(S)}{p_S}.
\]

\paragraph{Theorem (point-defect classification).}
Let $q_d:\{-1,1\}^d\to\{-1,1\}$ be negative at one cube point and
positive elsewhere, and write $\delta=2^{-d}$.  Then
\begin{align}
 I(q_d)&=2d\delta,\label{eq:I}\\
 U(q_d)&=(1-2\delta)^2
   \log_2\frac1{(1-2\delta)^2}
   +2d\delta(1-2\delta).\label{eq:U}
\end{align}
The first term in \eqref{eq:U} is interpreted as zero for $d=1$.  Moreover
\[
 H(q_d)<4I(q_d),\qquad U(q_d)<3I(q_d)                 \tag{3}
\]
for every $d\ge1$ (both inequalities are strict for every such $d$),
and
\[
 \frac{H(q_d)}{I(q_d)}\longrightarrow4,qquad
 \frac{U(q_d)}{I(q_d)}\longrightarrow1.              \tag{4}
\]

If $F=\prod_{j=1}^m q_{d_j}$ uses disjoint physical cubes, then
\[
 H(F)=\sum_jH(q_{d_j}),\quad I(F)=\sum_jI(q_{d_j}),
 \quad U(F)=\sum_jU(q_{d_j}).                          \tag{5}
\]
Consequently every finite tensor tower of point defects, with arbitrary
dimensions and multiplicities, satisfies
\[
 H(F)<4I(F),\qquad U(F)<3I(F).                         \tag{6}
\]
Thus no such tower is a counterfamily to FEI or to a bounded-$U/I$
sufficient route.

\paragraph{Proof of the exact one-defect formulas.}
After an input sign flip we may place the exceptional point at the all-one
vertex.  The squared Fourier coefficients are
\[
 p_\varnothing=(1-2\delta)^2,
 \qquad p_S=4\delta^2\quad(S\ne\varnothing).
\]
This gives \eqref{eq:I}.  For $S\ne\varnothing$,
\[
 \pi(S)=4\delta^2 2^{d-|S|},
 \qquad \log_2\frac{\pi(S)}{p_S}=d-|S|,
\]
while $\pi(\varnothing)=1$.  Since
\[
 \sum_{S\ne\varnothing}(d-|S|)=d2^{d-1}-d,
\]
the nonempty contribution to $U$ equals
\[
 4\delta^2(d2^{d-1}-d)=2d\delta(1-2\delta),
\]
which proves \eqref{eq:U}.

For $d\ge2$, set $r=2\delta$.  Direct expansion gives
\[
 H(q_d)=-(1-r)^2\log_2(1-r)^2+r(2-r)(2d-2),
 \qquad I(q_d)=dr.
\]
Using $-\log(1-r)\le r/(1-r)$, the first entropy term is at most
$2r(1-r)/\ln2<4r$.  On the other hand,
\[
 4dr-r(2-r)(2d-2)=4r+r^2(2d-2)>4r.
\]
Hence $H(q_d)<4I(q_d)$.  At $d=1$, $q_1$ is a signed dictator and
$H(q_1)=0$, $I(q_1)=1$.

For the potential, the same logarithmic inequality gives
\[
 (1-2\delta)^2\log_2\frac1{(1-2\delta)^2}
 \le \frac{4\delta(1-2\delta)}{\ln2}.
\]
Thus, for $d\ge2$,
\[
 \frac{U(q_d)}{I(q_d)}
 \le(1-2\delta)\left(1+\frac{2}{d\ln2}\right)
 <1+\frac1{\ln2}<3.
\]
The last inequality follows, for example, from $\ln2>1/2$.  The case
$d=1$ again has $U(q_1)=0$.  Formula \eqref{eq:U} and
$-\ln(1-2\delta)=2\delta+O(\delta^2)$ give the second limit in (4); the
first is the standard exact point-defect asymptotic and also follows from
the displayed formula for $H$.

For disjoint spectra, $p_{F}(S,T)=p_f(S)p_g(T)$ and
\[
 \pi_F(S,T)=\pi_f(S)\pi_g(T).
\]
Taking the logarithm and averaging proves additivity of $U$; the entropy
and influence identities are the usual product-spectrum laws.  This proves
(5)--(6).\qed

\paragraph{The A077 multiscale sequence.}
For the exact sequence in A077, take $d_k=2k$ and $N_k=2^{2k}$ copies at
each scale $3\le k\le K$.  With
$S_K=\sum_{k=3}^Kk$ and $\delta_k=2^{-2k}$, formulas (1)--(5) give
\[
 I(F_K)=4S_K,
\]
and the exact potential ledger
\[
 U(F_K)=\sum_{k=3}^K\left[
 \delta_k^{-1}(1-2\delta_k)^2
 \log_2\frac1{(1-2\delta_k)^2}
 +4k(1-2\delta_k)\right].                             \tag{7}
\]
Taylor's formula with a uniform geometric remainder yields
\[
 U(F_K)=4S_K+4(K-2)\log_2e+O(1),
 \qquad \frac{U(F_K)}{I(F_K)}\longrightarrow1.         \tag{8}
\]
The already derived exact entropy formula gives
\[
 H(F_K)=16S_K+O(K),
 \qquad \frac{H(F_K)}{I(F_K)}\longrightarrow4.         \tag{9}
\]
Finally the biased B09 gain satisfies $J_k/(r_kk)\to1$, while
$N_kr_k\to4$; hence its uniform transported aggregate is of the same order
as the physical influence, not larger by an unbounded factor.  The apparent
divergence occurs only after deleting the encoder label and boundary costs.

\paragraph{Scope.}
This theorem closes the entire disjoint tensor closure of Boolean point
defects, including arbitrary multiscale multiplicities.  It does not cover
overlapping nonlinear surgeries, selectors with non-singleton negative
cells, or a different gate whose Fourier upper-set law is not the one above.
Unbounded biased-helper cost is not an FEI counterexample, and full FEI
remains open.
